\chapter{INFORMATIONS G\'EN\'ERALES}

\section*{Pr\'esentation du projet}
\addcontentsline{toc}{section}{Pr\'esentation du projet}
AuditHub est une application web monoproduit destin\'ee \`a centraliser la documentation de missions de tests d'intrusion autoris\'es. Elle associe une interface React \`a une API JSON Django. Son r\^ole n'est pas d'ex\'ecuter Nmap, Metasploit, Burp Suite ou OWASP ZAP : les outils restent utilis\'es dans un laboratoire externe autoris\'e, puis leurs r\'esultats sont consign\'es sous forme textuelle dans l'application.

Le besoin trait\'e est la dispersion habituelle des informations d'un audit technique : autorisation, p\'erim\`etre IP, intervenants, m\'ethode, preuve, s\'ev\'erit\'e, responsable de correction, \`ech\'eance et re-test. AuditHub rassemble ces objets dans une cha\^ine tra\c{c}able et applique des contr\^oles d'acc\`es c\^ot\'e serveur.

\begin{table}[H]
\centering
\caption{Fiche d'identit\'e technique d'AuditHub}
\begin{tabularx}{\textwidth}{>{\bfseries}p{4.1cm}X}
\toprule
\rowcolor{AuditLight}Rubrique & Information v\'erifi\'ee\\
\midrule
Backend & Django 5.2.17, vues fonctionnelles, formulaires et ORM\\
Frontend & React 19.3.0, bundle ES2020 produit par esbuild 0.28.2\\
Persistance & SQLite, migration applicative jusqu'\`a \evidence{0002\_followup\_constraints}\\
Authentification & Sessions Django, cookie HttpOnly, jeton CSRF, limitation des \`echecs\\
R\^oles & Administrateur, Pentesteur, Auditeur\\
Tests & 33 tests Django et 3 parcours Playwright valid\'es le 15 octobre 2026\\
D\'eploiement document\'e & Waitress derri\`ere un reverse proxy TLS ; variables d'environnement obligatoires\\
\bottomrule
\end{tabularx}
\end{table}

\section*{1.3 Type d'audit}
\addcontentsline{toc}{section}{1.3 Type d'audit}
Le travail rapport\'e est une \textbf{revue technique et fonctionnelle du logiciel}, compl\'et\'ee par l'ex\'ecution de ses tests automatis\'es. Il ne s'agit pas d'un audit initial, de surveillance ou de renouvellement d'un SMSI au sens d'un organisme de certification. Cette qualification est essentielle : le d\'ep\^ot ne contient ni mandat de certification, ni D\'eclaration d'Applicabilit\'e organisationnelle, ni entretiens, ni preuves d'un environnement de production.

\section*{1.4 Dates et dur\'ee de l'audit}
\addcontentsline{toc}{section}{1.4 Dates et dur\'ee de l'audit}
L'\'etat final analys\'e et les validations pr\'esent\'ees dans ce rapport datent du \textbf{15 octobre 2026}. La date de cr\'eation du rapport de tests historique du projet est le 1\textsuperscript{er} octobre 2026, puis il a \`et\'e actualis\'e. La dur\'ee humaine totale de conception et de d\'eveloppement n'est pas enregistr\'ee dans le d\'ep\^ot ; elle ne peut donc pas \^etre chiffr\'ee sans supposition.

\section*{Probl\'ematique et besoins}
\addcontentsline{toc}{section}{Probl\'ematique et besoins}
La probl\'ematique peut \^etre formul\'ee ainsi : \emph{comment fournir une application simple, tra\c{c}able et contr\^ol\'ee pour documenter des tests de s\'ecurit\'e r\'ealis\'es hors plateforme, sans transformer le serveur web en outil offensif ?} La r\'eponse impl\'ement\'ee repose sur la s\'eparation entre ex\'ecution et documentation, un p\'erim\`etre explicite par mission, des r\^oles, la journalisation et un export JSON.
